%! Author = VORNAME NACHNAME
%! Date = TT.MM.JJ

% Preamble
\documentclass[12pt, a4paper, oneside]{article}

% Packages
\usepackage[utf8]{inputenc}
\usepackage[T1]{fontenc}
\usepackage[ngerman]{babel}
\usepackage{mathptmx}        % Times New Roman (Schriftart)
\usepackage{geometry}
\usepackage{graphicx}
\usepackage{setspace}
\usepackage{parskip}
\usepackage{lipsum}
\usepackage{titlesec}
\usepackage{fancyhdr}
\usepackage{hyperref}
\usepackage{xcolor}
\usepackage{tikz}
\usetikzlibrary{shapes.geometric, arrows.meta, fit, backgrounds, positioning, calc}
\usepackage{wrapfig}
\usepackage{float}
\usepackage{enumitem}
\usepackage{needspace}
\usepackage{acronym}
\usepackage{booktabs}            % Professionelle Tabellenlinien
\usepackage{tabularx}            % Flexible Tabellenbreite
\usepackage{colortbl}            % Farbige Tabellenzeilen
\usepackage{etoolbox}            % Makro-Werkzeuge (gappto etc.)
\usepackage[toc,title]{appendix} % Anhangsverzeichnis
\usepackage{listings}            % Quellcode-Listings (Anhang)

% Listing-Style für Quellcode-Auszüge im Anhang
\definecolor{lstbg}{rgb}{0.97,0.97,0.97}
\definecolor{lstkw}{rgb}{0.10,0.30,0.65}
\definecolor{lstcm}{rgb}{0.30,0.55,0.30}
\definecolor{lststr}{rgb}{0.65,0.15,0.15}
\lstdefinestyle{code}{
    backgroundcolor=\color{lstbg},
    basicstyle=\ttfamily\footnotesize,
    keywordstyle=\color{lstkw}\bfseries,
    commentstyle=\color{lstcm}\itshape,
    stringstyle=\color{lststr},
    numbers=left,
    numberstyle=\tiny\color{gray},
    numbersep=8pt,
    frame=single,
    rulecolor=\color{gray!40},
    breaklines=true,
    breakatwhitespace=true,
    columns=fullflexible,
    showstringspaces=false,
    tabsize=2,
    captionpos=b,
    aboveskip=8pt,
    belowskip=8pt,
}
\lstset{style=code}
\usepackage[
    backend=biber,
    style=verbose-ibid,   % volle bibliographische Angabe in der Fußnote, danach „ebd.“
    citetracker=true,
    sorting=nyt
]{biblatex}

% ============================================================
%  Deutsches Fußnoten-Zitierformat
%  Ziel: Nachname V.: Titel, Ort Jahr[, URL: <url> [aufgerufen TT.MM.JJJJ]], S. X
% ============================================================

% 1) Autorenformat → „Nachname V." (Nachname + Initiale mit Punkt)
\DeclareNameFormat{last-init}{%
  \namepartfamily
  \ifblank{\namepartgiveni}{}{\addspace\namepartgiveni}%
  \ifthenelse{\value{listcount}<\value{liststop}}
    {\addcomma\addspace}{}%
}
\DeclareNameAlias{author}{last-init}
\DeclareNameAlias{editor}{last-init}

% 2) Doppelpunkt nach Autorenname (in Zitationen UND Literaturverzeichnis)
%    nametitledelim wird in allen Kontexten auf Doppelpunkt gesetzt
\DeclareDelimFormat{nametitledelim}{\addcolon\addspace}
\DeclareDelimFormat[bib,biblist]{nametitledelim}{\addcolon\addspace}

% 3) Kein Verlag – nur Ort + Jahr: „Heidelberg 2002"
\renewbibmacro*{publisher+location+date}{%
  \printlist{location}%
  \iflistundef{location}{}{\setunit{\addspace}}%
  \usebibmacro{date}%
  \newunit
}

% 4) URL-Feld: „URL: <url>"
\DeclareFieldFormat{url}{URL\addcolon\addspace\url{#1}}

% 5) urldate: „[aufgerufen TT.MM.JJJJ]"
\DeclareFieldFormat{urldate}{%
  \mkbibbrackets{aufgerufen\addspace
    \iffieldundef{urlday}{}{\mkdatezeros{\thefield{urlday}}\adddot}%
    \iffieldundef{urlmonth}{}{\mkdatezeros{\thefield{urlmonth}}\adddot}%
    \thefield{urlyear}}%
}

% 7) In Zitationen: ISBN, ISSN, DOI und Auflage unterdrücken
\AtEveryCitekey{%
  \clearfield{isbn}%
  \clearfield{issn}%
  \clearfield{doi}%
  \clearfield{edition}%
  \clearfield{pagetotal}%
}

% ibidem-String
\DefineBibliographyStrings{ngerman}{%
    ibidem = {ebd\adddot}
}

\addbibresource{bibliography/literatur.bib}
\geometry{top=3cm, bottom=3cm, left=2.5cm, right=2cm}
\setlength{\headheight}{14.5pt}

\hypersetup{
    colorlinks=true,
    linkcolor=black,
    urlcolor=blue,
    citecolor=black
}

% Appendix-Namen auf Deutsch
\renewcommand{\appendixname}{Anhang}
\renewcommand{\appendixtocname}{Anhang}
\renewcommand{\appendixpagename}{Anhang}

% Header & Footer
\pagestyle{fancy}
\fancyhf{}
\rhead{VORNAME NACHNAME}           % TODO: Name anpassen
\lhead{MODULKÜRZEL -- Assignment}  % TODO: Modul anpassen
\cfoot{\thepage}

% Zähler um römische Seitenzahl am Ende fortzuführen
\newcounter{frontpages}

%-----------------------------------------------------------
% AI-VERZEICHNIS: Deklarations-System (analog zu .bib)
%  \aitool{Name}{Einsatzform}{Bemerkungen}
%  \printAIVerzeichnis  →  rendert die Tabelle (Nr. automatisch)
%-----------------------------------------------------------
\newcommand{\aitools}{}          % interner Datenpuffer
\newcounter{aitoolrow}           % Auto-Zähler (wird beim Rendern hochgezählt)
\newif\iffirstaitool             % Flag für ersten Eintrag
\firstaitooltrue

\newcommand{\aitool}[3]{%
    \iffirstaitool
        \global\firstaitoolfalse
    \else
        \gappto\aitools{\midrule}%
    \fi
    \gappto\aitools{\stepcounter{aitoolrow}\theaitoolrow & \textbf{#1} & #2 & #3 \\}%
}

\newcommand{\printAIVerzeichnis}{%
    \setcounter{aitoolrow}{0}%   % Zähler vor jedem Druck zurücksetzen
    \renewcommand{\arraystretch}{1.4}%
    \begin{tabularx}{\textwidth}{%
        >{\centering\arraybackslash}p{0.6cm}%
        p{3.4cm}%
        X%
        X%
    }%
    \toprule%
    \rowcolor{gray!20}%
    \textbf{Nr.} & \textbf{AI-Hilfsmittel} & \textbf{Einsatzform} & \textbf{Bemerkungen} \\%
    \midrule%
    \aitools%
    \bottomrule%
    \end{tabularx}%
    \renewcommand{\arraystretch}{1}%
}

% Datei laden (analog zu \addbibresource)
% AI-Verzeichnis -- Eintraege
% Syntax: \aitool{Hilfsmittel}{Einsatzform}{Bemerkungen}

% TODO: Genutzte KI-Tools eintragen, Beispiele:
% \aitool{GitHub Copilot}{Code-Generierung}{--}
% \aitool{ChatGPT}{Recherche und Formulierungshilfen}{--}


% Document
\begin{document}
\onehalfspacing

%-----------------------------------------------------------
% DECKBLATT (keine Seitenzahl)
%-----------------------------------------------------------
\pagenumbering{Roman}
\begin{titlepage}
    \thispagestyle{empty}
    \centering

    % Logo oben
    \vspace*{0.5cm}
    \includegraphics[width=0.45\textwidth]{assets/images/akad_logo}\\[1.5cm]

    % Trennlinie
    \textcolor{gray}{\rule{\linewidth}{0.4mm}}\\[0.6cm]

    % Titel
    {\fontsize{20}{24}\selectfont \textbf{TITEL DER ARBEIT}}\\[0.5cm]  % TODO

    % Untertitel
    {\large \textcolor{gray}{\textit{UNTERTITEL DER ARBEIT}}}\\[0.6cm]  % TODO

    \textcolor{gray}{\rule{\linewidth}{0.4mm}}\\[1.5cm]

    % Infobox
    \begin{tabular}{r@{\hspace{0.8cm}}l}
        \textbf{Verfasser}      & VORNAME NACHNAME \\[0.4cm]   % TODO
        \textbf{Matrikelnummer} & MATRIKELNUMMER \\[0.4cm]      % TODO
        \textbf{Modul}          & MODULKÜRZEL \\[0.4cm]         % TODO
        \textbf{Betreuer}       & Prof. Dr. BETREUER \\[0.4cm]  % TODO
        \textbf{Abgabedatum}    & TT. Monat JJJJ \\             % TODO
    \end{tabular}

    \vfill
\end{titlepage}

%-----------------------------------------------------------
% INHALTSVERZEICHNIS
%-----------------------------------------------------------
\clearpage
\pagestyle{empty}
\tableofcontents
\clearpage
\pagestyle{fancy}

%-----------------------------------------------------------
% ABBILDUNGSVERZEICHNIS - Römische Seitenzahl ab I
%-----------------------------------------------------------
\setcounter{page}{1}
\listoffigures
\addcontentsline{toc}{section}{Abbildungsverzeichnis}
\newpage

%-----------------------------------------------------------
% TABELLENVERZEICHNIS
%-----------------------------------------------------------
\listoftables
\addcontentsline{toc}{section}{Tabellenverzeichnis}
\newpage

%-----------------------------------------------------------
% ABKÜRZUNGSVERZEICHNIS
%-----------------------------------------------------------
\section*{Abkürzungsverzeichnis}
\addcontentsline{toc}{section}{Abkürzungsverzeichnis}
\label{sec:abkuerzungen}

\begin{acronym}[PLACEHOLDER]
    % TODO: Abkürzungen eintragen
    % Syntax: \acro{kürzel}[ANZEIGE]{Ausgeschriebener Begriff}
    % Beispiele:
    \acro{api}[API]{Application Programming Interface}
    \acro{sql}[SQL]{Structured Query Language}
\end{acronym}

% Römische Seitenzahl merken, bevor Hauptteil beginnt
\setcounter{frontpages}{\value{page}}
\newpage

%-----------------------------------------------------------
% INHALT - arabische Nummerierung ab 1
%-----------------------------------------------------------
\pagenumbering{arabic}
\section{Einleitung}
\label{sec:einleitung}
% TODO: Thema und Kontext der Arbeit einführen
\subsection{Hintergrund und Motivation}
\label{subsec:hintergrund}
% TODO: Hintergrund des Themas beschreiben
\subsection{Problemstellung}
\label{subsec:problemstellung}
% TODO: Konkrete Problemstellung formulieren
\subsection{Ziel der Arbeit}
\label{subsec:ziel}
% TODO: Ziel(e) der Arbeit benennen
% Folgende Fragestellungen sollen untersucht werden:
\begin{itemize}
    \item[--] Fragestellung 1
    \item[--] Fragestellung 2
    \item[--] Fragestellung 3
\end{itemize}
\subsection{Aufbau der Arbeit}
\label{subsec:aufbau}
% TODO: Kapitelstruktur kurz beschreiben

\newpage

\section{Grundlagen}
\label{sec:grundlagen}

% TODO: Theoretische Grundlagen des Themas erlaeutern

\subsection{Grundbegriffe}
\label{subsec:grundbegriffe}

% TODO: Zentrale Begriffe definieren

\subsection{Theoretischer Rahmen}
\label{subsec:theorie}

% TODO: Relevante Theorien / Konzepte beschreiben

\subsection{Stand der Forschung}
\label{subsec:forschungsstand}

% TODO: Vorhandene Literatur und Erkenntnisse zusammenfassen

\newpage

\section{Hauptteil}
\label{sec:hauptteil}
% TODO: Kernthema der Arbeit ausarbeiten
\subsection{Analyse}
\label{subsec:analyse}
% TODO: Analyse des Themas / der Fragestellung
\subsection{Ergebnisse}
\label{subsec:ergebnisse}
% TODO: Ergebnisse der Analyse darstellen
% Beispiel fuer eine Tabelle:
\begin{table}[H]
    \centering
    \renewcommand{\arraystretch}{1.3}
    \begin{tabularx}{\textwidth}{l X}
        \toprule
        \rowcolor{grayps -ef | grep Transmover}
        \textbf{Kriterium} & \textbf{Beschreibung} \
        \midrule
        Kriterium 1 & Beschreibung 1 \
        Kriterium 2 & Beschreibung 2 \
        \bottomrule
    \end{tabularx}
    \caption{TODO: Tabellenunterschrift}
    \label{tab:beispiel}
\end{table}
\subsection{Diskussion}
\label{subsec:diskussion}
% TODO: Ergebnisse kritisch einordnen und diskutieren

\newpage

\section{Fazit und Ausblick}
\label{sec:fazit}
\subsection{Zusammenfassung der wichtigsten Erkenntnisse}
\label{subsec:zusammenfassung}
% TODO: Zentrale Ergebnisse zusammenfassen und Fragestellungen aus der Einleitung beantworten
\subsection{Ausblick auf zukuenftige Entwicklungen}
\label{subsec:ausblick}
% TODO: Offene Fragen, Folgeforschung und Perspektiven benennen

\newpage

%-----------------------------------------------------------
% BACKMATTER - römische Nummerierung, fortführend
%-----------------------------------------------------------
\hypersetup{pageanchor=false}
\pagenumbering{Roman}
\setcounter{page}{\value{frontpages}}
\stepcounter{page}
\hypersetup{pageanchor=true}

%-----------------------------------------------------------
% LITERATURVERZEICHNIS
%-----------------------------------------------------------
\printbibliography[heading=bibintoc, title={Literaturverzeichnis}]
\newpage

%-----------------------------------------------------------
% AI-VERZEICHNIS
%-----------------------------------------------------------
\section*{AI-Verzeichnis}
\addcontentsline{toc}{section}{AI-Verzeichnis}

In der vorliegenden Arbeit wurden folgende KI-basierte Anwendungen und Werkzeuge genutzt:

\vspace{0.5cm}
\noindent
\printAIVerzeichnis
\newpage

%-----------------------------------------------------------
% ANHANGSVERZEICHNIS
%-----------------------------------------------------------
\addcontentsline{toc}{section}{Anhangsverzeichnis}
\listofappendices
\newpage

%-----------------------------------------------------------
% ANHANG
%-----------------------------------------------------------
\begin{appendices}
\section{Anhang A: Titel des Anhangs}
\label{sec:anhang-a}

% TODO: Anhang-Inhalt einfuegen (z.B. Code-Listings, Tabellen, Abbildungen)

% Beispiel fuer ein Code-Listing:
% \begin{lstlisting}[language=Python, caption={Beispiel-Listing}, label={lst:beispiel}]
% # Code hier einfuegen
% \end{lstlisting}

\end{appendices}

\newpage

%-----------------------------------------------------------
% EIDESSTATTLICHE ERKLÄRUNG (keine Seitenzahl im Verzeichnis)
%-----------------------------------------------------------
\thispagestyle{fancy}
\section*{Eidesstattliche Erklärung}

Ich versichere, dass ich die vorliegende Arbeit selbstständig verfasst, keine
anderen als die angegebenen Quellen und Hilfsmittel benutzt sowie alle wörtlich
oder sinngemäß übernommenen Stellen in der Arbeit gekennzeichnet habe. Alle
Teile meiner Arbeit, die wortwörtlich oder dem Sinn nach anderen Werken
entnommen sind, wurden unter Angabe der Quelle kenntlich gemacht. Gleiches gilt
auch für Zeichnungen, Skizzen, bildliche Darstellung sowie für Quellen aus dem
Internet, dazu zählen auch KI-basierte Anwendungen oder Werkzeuge. Die Arbeit
wurde in gleicher oder ähnlicher Form noch nicht als Prüfungsleistung eingereicht.

\vspace{2.5cm}

\noindent
\begin{tabular}{p{6cm} p{2cm} p{6cm}}
    \hrulefill & & \hrulefill \\[0.3cm]
    Ort, Datum & & Unterschrift
\end{tabular}

\end{document}
